\begin{example}[Ticker Engine]\label{ex:ticker-engine}%
    Let us define a simple ticker engine that processes increment and count
    messages. This example demonstrates basic state management and actor
    spawning with a continuation pattern.
    
    First, we define the message interface for the ticker engine:
    
    \begin{equation}\label{eq:ticker-message-interface}
    \begin{aligned}
    \typedef%
    {
      \msgtype{\mathsf{ticker}}
    }{    &  \msg{Increment}() \\
    \mid &  \msg{Count}() \\
    \mid &  \msg{CountResponse}(\N)
    }\,,
    \end{aligned}
    \end{equation}
    
    where:
    \begin{itemize}
        \item $\msg{Increment}()$ increments the internal counter
        \item $\msg{Count}()$ requests the current counter value
        \item $\msg{CountResponse}(value)$ returns the current counter value
    \end{itemize}
    
    The configuration type for the ticker engine:
    
    \begin{equation}\label{eq:ticker-configuration}
    \begin{aligned}
    \typedef{%
      \ctype{C}_{\mathsf{ticker}}
    }{    &  \ccon{Standard}(\N) 
    }\,.
    \end{aligned}
    \end{equation}
    
    where the natural number parameter specifies the starting value for the
    ticker.
    
    The environment type is quite simple and just tracks the ticker's count:
    
    \begin{equation}\label{eq:ticker-environment}
    \begin{aligned}
    \typedef{%
      \ctype{L}_{\mathsf{ticker}}
    }{    &  \ccon{TickerState}(\N)
    }\,,
    \end{aligned}
    \end{equation}
    
    where:
    \begin{itemize}
        \item $\ccon{TickerState}$ tracks the current count value
    \end{itemize}
    
    The guard for processing increment messages:
    
    \begin{equation}\label{eq:ticker-increment-guard}
    \begin{aligned}
    &\typing{
      g_{\mathsf{increment}}
    }{
      \In{\mathsf{ticker}}
        \to
      \Option{\Unit}
    }\,. \\
    &\defequal{
      g_{\mathsf{increment}}
    }{
      \lambda \langle m, \ignore, \ignore \rangle .
      \mathsf{case}\ m\ \mathsf{of}\ \{ \\
    &  \msg{Increment}() \Rightarrow \some{()}; \\
    &  \ignore \Rightarrow \none \\
    &\}
    }\,.
    \end{aligned}
    \end{equation}
    
    The corresponding action increments the counter and spawns a new ticker when
    reaching 100:
    
    \begin{equation}\label{eq:ticker-increment-action}
    \begin{aligned}
    &\typing{
      a_{\mathsf{increment}}
    }{
      \Unit \times \In{\mathsf{ticker}}
        \to \Out{\mathsf{ticker}}
    }\,. \\
    &\defequal{
      a_{\mathsf{increment}}
    }{
      \lambda \langle \ignore, \langle \mathsf{sender}, \ignore, \langle \mathsf{count} \rangle \rangle \rangle . \\
    &  \chain{ \\
    &    \mathsf{let}\ newCount = \mathsf{count} + 1 \\
    &    \update{\langle newCount \rangle} \\
    &    \mathsf{if}\ newCount \div 100 = 0 \land newCount > 0\ \mathsf{then} \\
    &      \chain{ \\
    &        \spawn{ \\
    &          \configval{%
                \some{\self}%
              }{%
                \mathsf{ticker}%
              }{%
                \ccon{Standard}(newCount)%
              }%
            }{%
              \localenvval{%
                \ccon{TickerState}(newCount)%
              }{%
                [(\self, \ignore)]%
              }%
            }%
          }{ \\
    &        \send{\mathsf{sender}}{\msg{CountResponse}(newCount)} \\
    &        } \\
    &    } \\
    &  \}
    }\,.
    \end{aligned}
    \end{equation}
    
    The guard for processing count messages:
    
    \begin{equation}\label{eq:ticker-count-guard}
    \begin{aligned}
    &\typing{
      g_{\mathsf{count}}
    }{
      \In{\mathsf{ticker}}
        \to
      \Option{\Unit}
    }\,. \\
    &\defequal{
      g_{\mathsf{count}}
    }{
      \lambda \langle m, \ignore, \ignore \rangle .
      \mathsf{case}\ m\ \mathsf{of}\ \{ \\
    &  \msg{Count}() \Rightarrow \some{()}; \\
    &  \ignore \Rightarrow \none \\
    &\}
    }\,.
    \end{aligned}
    \end{equation}
    
    The corresponding action returns the current count:
    
    \begin{equation}\label{eq:ticker-count-action}
    \begin{aligned}
    &\typing{
      a_{\mathsf{count}}
    }{
      \Unit \times \In{\mathsf{ticker}}
        \to \Out{\mathsf{ticker}}
    }\,. \\
    &\defequal{
      a_{\mathsf{count}}
    }{
      \lambda \langle \ignore, \langle \mathsf{sender}, \ignore, \langle \mathsf{count} \rangle \rangle \rangle . \\
    &  \send{\mathsf{sender}}{\msg{CountResponse}(\mathsf{count})}
    }\,.
    \end{aligned}
    \end{equation}
    
    The initialization action for the ticker:
    
    \begin{equation}\label{eq:ticker-init-action}
    \begin{aligned}
    &\typing{
      a_{\mathsf{init}}
    }{
      \In{\mathsf{ticker}}
        \to \Out{\mathsf{ticker}}
    }\,. \\
    &\defequal{
      a_{\mathsf{init}}
    }{
      \lambda \langle \ignore, \ccon{Standard}(startValue), \langle \mathsf{count} \rangle \rangle . \\
    &  \update{\langle startValue \rangle}
    }\,.
    \end{aligned}
    \end{equation}
    
    Finally, we define the complete behaviour of the ticker engine:
    
    \begin{equation}\label{eq:ticker-behaviour}
    \begin{aligned}
    \defequal{
      \typing{\behaviour{\mathsf{ticker}}}{\behaviourtype{\mathsf{ticker}}}
    }{
      [ a_{\mathsf{init}}
      , g_{\mathsf{increment}} \odot a_{\mathsf{increment}}
      , g_{\mathsf{count}} \odot a_{\mathsf{count}}
      ]
    }\,.
    \end{aligned}
    \end{equation}
    
    This example demonstrates a simple stateful actor that can spawn
    continuations of itself. When the count reaches multiples of 100, the ticker
    spawns a new ticker to continue counting from the current value. This
    pattern could be useful for processes that need to be periodically refreshed
    or migrated to maintain system health.
\end{example} 